\chapter{Esteiras de CI/CD Multi-Ambiente para Produtos AWS Banc\'arios}
\label{chap:cicd_multiambiente_bancario}

% ---------------------------------------------------------------
\section{Escopo do cap\'itulo}
% ---------------------------------------------------------------

Este cap\'itulo descreve a organiza\c{c}\~ao e a opera\c{c}\~ao das
esteiras de integra\c{c}\~ao cont\'inua\index{integração contínua|textbf}\index{CI/CD|textbf} e entrega cont\'inua
(\textit{Continuous Integration/Continuous Delivery} --- CI/CD) adotadas
nos produtos AWS\index{AWS} do projeto bancário. A estratégia baseia-se em
tr\^es ambientes isolados --- \texttt{dev}\index{dev (ambiente)|textbf},
\texttt{homol}\index{homol (ambiente)|textbf} e \texttt{prod}\index{prod (ambiente)|textbf} ---
cada qual mapeado para uma conta AWS\index{conta AWS separada} independente,
garantindo isolamento de recursos, permiss\~oes e dados entre est\'agios do
ciclo de vida do software. A infraestrutura de cada produto \'e integralmente
descrita em m\'odulos \textit{Terraform}\index{Terraform} versionados, e os
fluxos de automa\c{c}\~ao s\~ao implementados como
\textit{Git Workflows}\index{Git Workflow|textbf}\index{GitHub Actions|textbf}
(\textit{GitHub Actions} ou \textit{GitLab CI}\index{GitLab CI}).

% ---------------------------------------------------------------
\section{Objetivos de aprendizagem}
% ---------------------------------------------------------------

Ao concluir este cap\'itulo, o leitor ser\'a capaz de:

\begin{itemize}[leftmargin=*]
  \item Compreender a arquitetura de tr\^es ambientes AWS separados e
        as implica\c{c}\~oes de isolamento de conta para sistemas banc\'arios;
  \item Estruturar esteiras de CI/CD independentes por produto AWS,
        com est\'agios de build, valida\c{c}\~ao, homologa\c{c}\~ao e
        produ\c{c}\~ao claramente definidos;
  \item Configurar \textit{workspaces} Terraform\index{Terraform!workspace|textbf}
        mapeados a ambientes e \textit{state backends}\index{state backend|textbf}
        segregados por conta;
  \item Implementar controles de promo\c{c}\~ao entre ambientes com
        aprova\c{c}\~ao humana obrigat\'oria antes de \texttt{homol}
        e \texttt{prod}.
\end{itemize}

% ---------------------------------------------------------------
\section{Conceitos te\'oricos}
% ---------------------------------------------------------------

\subsection{Isolamento de contas AWS por ambiente}

O modelo de tr\^es contas AWS separadas --- uma por ambiente ---
constitui pr\'atica recomendada pelo \textit{AWS Well-Architected
Framework}\index{AWS Well-Architected Framework} para sistemas
financeiros regulados. O isolamento por conta oferece:

\begin{itemize}[leftmargin=*]
  \item \textbf{Blast radius controlado:} falhas, vazamentos de
        credenciais ou configura\c{c}\~oes equivocadas em
        \texttt{dev} n\~ao afetam \texttt{prod};
  \item \textbf{Pol\'iticas IAM\index{IAM} independentes:} cada conta
        define seus pr\'oprios limites de permiss\~ao sem interfer\^encia
        cruzada;
  \item \textbf{Segregação de dados:} dados reais de clientes presentes
        apenas em \texttt{prod}; \texttt{dev} e \texttt{homol} operam
        exclusivamente com dados sint\'eticos ou anonimizados.
\end{itemize}

O fluxo de promo\c{c}\~ao entre ambientes \'e descrito na subse\c{c}\~ao seguinte.

\subsection{Esteira por produto}

Cada produto AWS --- Lambda\index{Lambda (AWS)}, API Gateway\index{API Gateway},
DynamoDB\index{DynamoDB}, SQS\index{SQS}, OpenSearch\index{OpenSearch},
Neptune\index{Neptune (AWS)} --- mant\'em seu pr\'oprio reposit\'orio Git
com conjunto de arquivos de \textit{workflow} independente. Essa separação
previne acoplamento entre produtos e permite que equipes distintas evoluam
cadências de deploy sem bloquear umas às outras.

\subsection{Fluxo de promoção entre ambientes}

O ciclo de promo\c{c}\~ao segue o padr\~ao:

\begin{enumerate}[leftmargin=*]
  \item \textbf{Feature branch $\to$ \texttt{dev}:} merge disparado por
        \textit{pull request}\index{pull request} aprovado; pipeline executa
        lint, testes unit\'arios e \texttt{terraform plan}.
  \item \textbf{\texttt{dev} $\to$ \texttt{homol}:} merge na branch
        \texttt{main}; exige aprova\c{c}\~ao do respons\'avel de QA e
        pipeline verde.
  \item \textbf{\texttt{homol} $\to$ \texttt{prod}:} \textit{tag} de
        release sem\'antica\index{versionamento semântico|textbf}
        (\texttt{vMAJOR.MINOR.PATCH}); exige aprova\c{c}\~ao de dois
        revisores autorizados e janela de manuten\c{c}\~ao registrada.
\end{enumerate}

% ---------------------------------------------------------------
\section{Implementa\c{c}\~ao orientada por passos}
% ---------------------------------------------------------------

\subsection{Passo 1 --- Estrutura de reposit\'orio por produto}

A organiza\c{c}\~ao recomendada para o reposit\'orio de um produto AWS \'e:

\begin{lstlisting}[language=bash,
  caption={Estrutura de reposit\'orio de produto AWS banc\'ario.},
  label={lst:repo_estrutura}]
produto-lambda-segmentacao/
  .github/
    workflows/
      ci.yml          # lint, testes, plan -- toda PR
      deploy-dev.yml  # deploy automatico em dev
      deploy-homol.yml# deploy com aprovacao QA
      deploy-prod.yml # deploy com aprovacao dupla + tag
  infraestrutura/
    main.tf
    variables.tf
    outputs.tf
    backend.tf        # configuracao S3 + DynamoDB lock
    envs/
      dev.tfvars
      homol.tfvars
      prod.tfvars     # sem segredos -- apenas referencias
  src/
    lambda_handler.py
    tests/
      test_unit.py
      test_integration.py
  requirements.txt
  requirements-dev.txt
\end{lstlisting}

\subsection{Passo 2 --- Configurar \textit{backend} Terraform por ambiente}

Cada ambiente utiliza um \textit{bucket} S3\index{S3 (AWS)} e uma tabela
DynamoDB dedicados para armazenar e bloquear o \textit{state}:

\begin{lstlisting}[language=bash,
  caption={Backend Terraform com lock por conta AWS.},
  label={lst:terraform_backend}]
# backend.tf
terraform {
  backend "s3" {
    bucket         = "tfstate-${var.ambiente}-123456789"
    key            = "produto-lambda-segmentacao/terraform.tfstate"
    region         = "us-east-1"
    dynamodb_table = "tfstate-lock-${var.ambiente}"
    encrypt        = true
  }
}
\end{lstlisting}

O par\^ametro \texttt{encrypt = true} garante criptografia do
\textit{state file}\index{state file (Terraform)|textbf} em repouso,
exig\^encia obrigat\'oria para reposit\'orios banc\'arios que possam
conter ARNs, IDs de recursos ou refer\^encias a segredos.

\subsection{Passo 3 --- Workflow de CI (pull request)}

O workflow de CI \'e disparado em toda \textit{pull request} aberta
contra qualquer branch protegida:

\begin{lstlisting}[language=bash,
  caption={GitHub Actions: workflow de CI para produto AWS.},
  label={lst:ci_workflow}]
# .github/workflows/ci.yml
name: CI

on:
  pull_request:
    branches: [main, develop]

jobs:
  quality:
    runs-on: ubuntu-latest
    steps:
      - uses: actions/checkout@v4

      - name: Setup Python
        uses: actions/setup-python@v5
        with:
          python-version: "3.11"

      - name: Instalar dependencias
        run: pip install -r requirements-dev.txt

      - name: Lint (flake8 + black)
        run: |
          flake8 src/ --max-line-length=100
          black --check src/

      - name: Testes unitarios
        run: pytest src/tests/test_unit.py --cov=src --cov-fail-under=70

      - name: Seguranca estatica (bandit)
        run: bandit -r src/ -ll

      - name: Terraform validate
        uses: hashicorp/setup-terraform@v3
        with:
          terraform_version: "1.7.0"
      - run: |
          cd infraestrutura
          terraform init -backend=false
          terraform fmt -check -recursive
          terraform validate
\end{lstlisting}

\subsection{Passo 4 --- Workflow de deploy com promo\c{c}\~ao controlada}

O deploy em \texttt{homol} e \texttt{prod} utiliza o mecanismo de
\textit{environment protection rules}\index{environment protection rules}
do GitHub Actions, que bloqueia a execu\c{c}\~ao at\'e que revisores
autorizados aprovem:

\begin{lstlisting}[language=bash,
  caption={GitHub Actions: deploy em homol com aprovacao.},
  label={lst:deploy_homol}]
# .github/workflows/deploy-homol.yml
name: Deploy Homol

on:
  push:
    branches: [main]

jobs:
  deploy:
    runs-on: ubuntu-latest
    environment: homologacao   # protection rule: 1 aprovador QA
    permissions:
      id-token: write          # OIDC para assume role sem chave estatica
      contents: read

    steps:
      - uses: actions/checkout@v4

      - name: Configurar credenciais AWS via OIDC
        uses: aws-actions/configure-aws-credentials@v4
        with:
          role-to-assume: arn:aws:iam::HOMOL_ACCOUNT_ID:role/GitHubActionsRole
          aws-region: us-east-1

      - name: Terraform apply (homol)
        run: |
          cd infraestrutura
          terraform init
          terraform workspace select homol
          terraform apply -var-file=envs/homol.tfvars -auto-approve
\end{lstlisting}

A autentica\c{c}\~ao via \textit{OIDC}\index{OIDC|textbf}\index{OpenID Connect|see{OIDC}}
(\textit{OpenID Connect}) elimina a necessidade de chaves de acesso
est\'aticas (\texttt{AWS\_ACCESS\_KEY\_ID} / \texttt{AWS\_SECRET\_ACCESS\_KEY})
armazenadas como \textit{secrets} do reposit\'orio, reduzindo
significativamente a superf\'icie de ataque\index{superfície de ataque}.

% ---------------------------------------------------------------
\section{Plano de testes do cap\'itulo}
% ---------------------------------------------------------------

\subsection{Teste funcional}

\begin{itemize}[leftmargin=*]
  \item \textbf{TC-01:} Abrir uma PR de feature, introduzir um erro de
        lint intencional e verificar que o job \texttt{quality} falha
        antes do deploy.
  \item \textbf{TC-02:} Fazer merge na branch \texttt{main} sem
        aprova\c{c}\~ao do QA e verificar que o job \texttt{deploy-homol}
        fica bloqueado aguardando aprova\c{c}\~ao.
  \item \textbf{TC-03:} Criar uma tag \texttt{v1.0.0} sem os dois
        revisores autorizados e verificar que o deploy em \texttt{prod}
        n\~ao prossegue.
\end{itemize}

\subsection{Teste de desempenho}

\begin{itemize}[leftmargin=*]
  \item \textbf{TP-01:} Medir o tempo total da pipeline de CI (checkout
        a valida\c{c}\~ao Terraform); crit\'erio: $\leq 8$\,min.
  \item \textbf{TP-02:} Medir o tempo de \texttt{terraform apply} em
        \texttt{homol} para o produto Lambda de segmenta\c{c}\~ao;
        crit\'erio: $\leq 5$\,min.
\end{itemize}

\subsection{Teste de qualidade}

\begin{itemize}[leftmargin=*]
  \item \textbf{TQ-01:} Executar \texttt{checkov} contra os m\'odulos
        Terraform e verificar zero falhas de n\'ivel CRITICAL.
  \item \textbf{TQ-02:} Validar que o \textit{state file} de
        \texttt{prod} est\'a criptografado e o acesso ao bucket S3
        \'e restrito apenas \`a role de CI/CD.
\end{itemize}

% ---------------------------------------------------------------
\section{Crit\'erios de aceite}
% ---------------------------------------------------------------

\begin{itemize}[leftmargin=*]
  \item TC-01 a TC-03 aprovados sem interven\c{c}\~ao manual.
  \item TP-01: pipeline de CI $\leq 8$\,min evidenciada com
        \textit{timestamps} do GitHub Actions.
  \item TQ-01: zero falhas CRITICAL no \texttt{checkov}.
  \item Nenhuma chave de acesso est\'atica AWS presente nos
        \textit{secrets} do reposit\'orio; autentica\c{c}\~ao
        exclusivamente via OIDC.
\end{itemize}

% ---------------------------------------------------------------
\section{Espa\c{c}o para expans\~ao iterativa}
% ---------------------------------------------------------------

\begin{itemize}[leftmargin=*]
  \item \textbf{v15.1} --- Implementar \textit{drift detection}\index{drift detection}:
        job agendado que executa \texttt{terraform plan} diariamente em
        \texttt{prod} e alerta via SNS\index{SNS (AWS)} caso detecte desvio
        entre o \textit{state} e a infraestrutura real.
  \item \textbf{v15.2} --- Adotar \textit{GitOps}\index{GitOps} com
        \textit{ArgoCD} ou \textit{Flux} para sincroniza\c{c}\~ao
        declarativa de manifests Kubernetes, se a plataforma evoluir
        para EKS\index{EKS (AWS)}.
  \item \textbf{Lacuna} --- Investigar o impacto de \textit{feature flags}
        no ciclo de promo\c{c}\~ao, permitindo ativar funcionalidades
        seletivamente por conta AWS sem novo deploy.
\end{itemize}
